\documentclass[12pt,a4paper]{article}
\usepackage[utf8]{inputenc}
\usepackage[spanish]{babel}
\usepackage{geometry}
\usepackage{listings}
\usepackage{xcolor}
\usepackage{parskip}
\usepackage{hyperref}

\geometry{margin=2.5cm}

\lstset{
  basicstyle=\ttfamily\small,
  keywordstyle=\color{blue}\bfseries,
  commentstyle=\color{gray},
  breaklines=true, frame=single, language=Java,
  numbers=left, numberstyle=\tiny\color{gray}
}

\title{\textbf{Puzzle con DFS Recursivo}\\
  \large Documentación Técnica}
\author{Inteligencia Artificial}
\date{}

\begin{document}
\maketitle
\tableofcontents
\newpage

\section{Descripción del Proyecto}
Visualizador comparativo de dos algoritmos de búsqueda aplicados a un puzzle: \textbf{DFS Recursivo} y \textbf{BFS}. Permite ver animadamente cómo cada algoritmo explora el espacio de estados hasta encontrar la solución.

\section{DFS Recursivo}

\subsection{Concepto}
El DFS Recursivo explora tan profundo como sea posible por cada rama antes de retroceder (backtrack). Su implementación recursiva delega el manejo de la pila al sistema operativo.

\subsection{Implementación}
\begin{lstlisting}
Nodo_DFS resolverDFS(Nodo_DFS nodo,
                     List<Nodo_DFS> visitados) {
    if (esMeta(nodo.getEstado())) return nodo;

    for (int[] estadoHijo : generarHijos(
                               nodo.getEstado())) {
        if (!estaEnLista(visitados, estadoHijo)) {
            Nodo_DFS hijo = new Nodo_DFS(estadoHijo);
            hijo.setPadre(nodo);
            visitados.add(hijo);

            Nodo_DFS resultado =
                resolverDFS(hijo, visitados);
            if (resultado != null) return resultado;
            // si retorna null, hace backtrack
        }
    }
    return null; // no encontró solución por esta rama
}
\end{lstlisting}

\subsection{Ventajas y desventajas}
\begin{itemize}
  \item \textbf{Ventaja}: código simple y elegante, menos memoria que BFS.
  \item \textbf{Desventaja}: no garantiza solución óptima; riesgo de \texttt{StackOverflowError} en problemas profundos.
\end{itemize}

\section{BFS para Comparación}

Se implementa BFS en la misma GUI para comparar visualmente cuántos pasos tarda cada algoritmo:

\begin{lstlisting}
Nodo_BFS resolverBFS(int[] estadoInicial) {
    Queue<Nodo_BFS> cola = new LinkedList<>();
    List<Nodo_BFS> visitados = new ArrayList<>();
    Nodo_BFS raiz = new Nodo_BFS(estadoInicial);
    cola.add(raiz);
    visitados.add(raiz);

    while (!cola.isEmpty()) {
        Nodo_BFS actual = cola.poll();
        if (esMeta(actual.getEstado())) return actual;

        for (int[] hijo : generarHijos(
                             actual.getEstado())) {
            if (!estaVisitado(visitados, hijo)) {
                Nodo_BFS nHijo = new Nodo_BFS(hijo);
                nHijo.setPadre(actual);
                visitados.add(nHijo);
                cola.add(nHijo);
            }
        }
    }
    return null;
}
\end{lstlisting}

\section{Animación de la Solución}

Una vez encontrada la solución, se reconstruye el camino y se anima:

\begin{lstlisting}
void animarSolucion(List<int[]> camino) {
    int[] paso = {0};
    Timer timer = new Timer(500, null);
    timer.addActionListener(e -> {
        if (paso[0] < camino.size()) {
            estadoActual = camino.get(paso[0]);
            repaint();
            paso[0]++;
        } else {
            ((Timer)e.getSource()).stop();
        }
    });
    timer.start();
}
\end{lstlisting}

\section{Interfaz Gráfica}

\begin{itemize}
  \item Visualización del puzzle con fichas numeradas.
  \item Botón \textbf{Auto DFS}: ejecuta y anima la solución con DFS recursivo.
  \item Botón \textbf{Auto BFS}: ejecuta y anima la solución con BFS.
  \item Contador de pasos para comparar ambos algoritmos.
\end{itemize}

\end{document}
